\uuid{gplS}
\titre{Équation droite/plan}
\auteur{bodin ; hanani ; mzari}
\organisation{exo7}
\licence{CC BY-NC-SA 4.0}
\datecreate{2020-01-20}
\niveau{L1}
\module{Géométrie}
\chapitre{Géométrie de l'espace}
\sousChapitre{Plans}
\difficulte{1}
\theme{Plans, équation droite, plan}
\exo7qcmid{lille-1.geometrie-espace.006}

\contenu{
\texte{
Pour ces questions, l'espace est muni d'un repère orthonormé direct  $(O, \vec{i}, \vec{j}, \vec{k})$.
}

\begin{question}
\qclassification{141.03}
\qtags{motcle=équation droite/plan}
Soit $P$ le plan passant par les points $A(0,1,0)$, $B(1,-1,0)$ et $C(0,1,1)$. Quelles sont les assertions vraies ?
\begin{answers}

    \good {Une représentation paramétrique de $P$ est :
 $$\left\{\begin{array}{ccl}x&=&s\\y&=&1-2s\\ z&=&t,\quad (t,s\in \Rr)\end{array}\right.$$}

    \bad {Une représentation paramétrique de $P$ est :
 $$\left\{\begin{array}{ccl}x&=&t\\y&=&s\\ z&=&1+2s,\quad (t,s\in \Rr)\end{array}\right.$$}

   \bad{Une équation cartésienne de $P$ est $2x+z=1$.}

    \good{Une équation cartésienne de $P$ est $2x+y=1$.}

\end{answers}
\begin{explanations}
$P$ est le plan passant par $A$ et dirigé par les vecteurs $\overrightarrow{AB}$ et $\overrightarrow{AC}$.
\end{explanations}
\end{question}
}
